% GENERATED FILE. Edit canonical lessons, then run npm run generate:books.
% canonical-source-hash: fnv1a64:19de944341641039
% canonical-lessons: GE-C321-erfinden, GE-C321-erfindung, GE-C321-gruenden, GE-C321-entstehen, GE-C321-zurueckkehren

\chapter{Neues entsteht}
\label{ch:ge-neues-entsteht}
\hlchaptermodality{321}

\begin{chapteropening}
\textbf{By the end of this chapter:} \emph{I can tell how things began: to invent, an invention, to found, to come into being and to return.}

Complete the last lesson of chapter 321: I can tell how things began: to invent, an invention, to found, to come into being and to return.
\end{chapteropening}

\section[erfinden]{erfinden — to invent; to make up}

\label{lesson:GE-C321-erfinden}



\begin{quote}
Before the new one: say the German for the ancestor, then the German for the downfall.
\end{quote}

\subsection*{You'll want to know: erfinden}
\textbf{erfinden} — \textquotedblleft{}to invent; to make up\textquotedblright{}.

\begin{quote}
Sie erfand eine neue Maschine.
\end{quote}

\begin{quote}
\emph{She invented a new machine.}
\end{quote}

Strong: \emph{erfand}, \emph{hat erfunden}. Also to make up a story.

\begin{cousinweb}
\emph{erfinden} is built on \emph{finden} \textquotedblleft{}to find\textquotedblright{}, a cousin of English \emph{find} (EtymDB), with \emph{er-}: to find out something new.
\end{cousinweb}

\begin{grammarlens}[title={What you've built}]
One more word for telling, explaining and describing at B1.
\end{grammarlens}

\subsection*{Your turn}
\begin{itemize}
\raggedright
  \item \emph{Say it:} \emph{erfinden}
  \item \emph{Say it:} \emph{erfinden}, once more
  \item \emph{Say it:} a sentence of your own with \emph{erfinden}
  \item \emph{Recall:} say the German for the ancestor, then the German for the downfall, then say \emph{erfinden} again
\end{itemize}

\begin{tcolorbox}[breakable,colback=teal!4,colframe=teal!35!black,title={Before you move on}]
What does erfinden mean? (\textquotedblleft{}To invent; to make up\textquotedblright{}.) Say it once more.
\end{tcolorbox}
\section[die Erfindung]{die Erfindung — the invention}

\label{lesson:GE-C321-erfindung}



\begin{quote}
Before the new one: say the German for the downfall, then the German for to invent.
\end{quote}

\subsection*{You'll want to know: die Erfindung}
\textbf{die Erfindung} — \textquotedblleft{}the invention\textquotedblright{}.

\begin{quote}
Das Rad war eine wichtige Erfindung.
\end{quote}

\begin{quote}
\emph{The wheel was an important invention.}
\end{quote}

Plural: \emph{die Erfindungen}.

\begin{cousinweb}
\emph{Erfindung} is built on \emph{erfinden} \textquotedblleft{}to invent\textquotedblright{}, from \emph{finden} \textquotedblleft{}to find\textquotedblright{}, a cousin of English \emph{find} (EtymDB).
\end{cousinweb}

\begin{grammarlens}[title={What you've built}]
One more word for telling, explaining and describing at B1.
\end{grammarlens}

\subsection*{Your turn}
\begin{itemize}
\raggedright
  \item \emph{Say it:} \emph{die Erfindung}
  \item \emph{Say it:} \emph{die Erfindung}, once more
  \item \emph{Say it:} a sentence of your own with \emph{die Erfindung}
  \item \emph{Recall:} say the German for the downfall, then the German for to invent, then say \emph{die Erfindung} again
\end{itemize}

\begin{tcolorbox}[breakable,colback=teal!4,colframe=teal!35!black,title={Before you move on}]
What does die Erfindung mean? (\textquotedblleft{}The invention\textquotedblright{}.) Say it once more.
\end{tcolorbox}
\section[gründen]{gründen — to found, to set up}

\label{lesson:GE-C321-gruenden}



\begin{quote}
Before the new one: say the German for to invent, then the German for the invention.
\end{quote}

\subsection*{You'll want to know: gründen}
\textbf{gründen} — \textquotedblleft{}to found, to set up\textquotedblright{}.

\begin{quote}
Mein Großvater gründete die Firma vor fünfzig Jahren.
\end{quote}

\begin{quote}
\emph{My grandfather founded the company fifty years ago.}
\end{quote}

\emph{Die Gründung} — the founding.

\begin{cousinweb}
\emph{gründen} is built on \emph{Grund} \textquotedblleft{}ground\textquotedblright{}, an inherited Germanic word from Proto-Germanic \emph{grunduz} (EtymDB), a cousin of English \emph{ground}: to lay the ground.
\end{cousinweb}

\begin{grammarlens}[title={What you've built}]
One more word for telling, explaining and describing at B1.
\end{grammarlens}

\subsection*{Your turn}
\begin{itemize}
\raggedright
  \item \emph{Say it:} \emph{gründen}
  \item \emph{Say it:} \emph{gründen}, once more
  \item \emph{Say it:} a sentence of your own with \emph{gründen}
  \item \emph{Recall:} say the German for to invent, then the German for the invention, then say \emph{gründen} again
\end{itemize}

\begin{tcolorbox}[breakable,colback=teal!4,colframe=teal!35!black,title={Before you move on}]
What does gründen mean? (\textquotedblleft{}To found, to set up\textquotedblright{}.) Say it once more.
\end{tcolorbox}
\section[entstehen]{entstehen — to come into being, to arise}

\label{lesson:GE-C321-entstehen}



\begin{quote}
Before the new one: say the German for the invention, then the German for to found.
\end{quote}

\subsection*{You'll want to know: entstehen}
\textbf{entstehen} — \textquotedblleft{}to come into being, to arise\textquotedblright{}.

\begin{quote}
Die Stadt entstand im Mittelalter.
\end{quote}

\begin{quote}
\emph{The town came into being in the Middle Ages.}
\end{quote}

Strong, with \emph{sein}: \emph{entstand}, \emph{ist entstanden}.

\begin{cousinweb}
\emph{entstehen} is built on \emph{stehen} \textquotedblleft{}to stand\textquotedblright{}, from Proto-Germanic \emph{stāną} (EtymDB), with \emph{ent-}: to come to stand.
\end{cousinweb}

\begin{grammarlens}[title={What you've built}]
One more word for telling, explaining and describing at B1.
\end{grammarlens}

\subsection*{Your turn}
\begin{itemize}
\raggedright
  \item \emph{Say it:} \emph{entstehen}
  \item \emph{Say it:} \emph{entstehen}, once more
  \item \emph{Say it:} a sentence of your own with \emph{entstehen}
  \item \emph{Recall:} say the German for the invention, then the German for to found, then say \emph{entstehen} again
\end{itemize}

\begin{tcolorbox}[breakable,colback=teal!4,colframe=teal!35!black,title={Before you move on}]
What does entstehen mean? (\textquotedblleft{}To come into being, to arise\textquotedblright{}.) Say it once more.
\end{tcolorbox}
\section[zurückkehren]{zurückkehren — to return, to come back}

\label{lesson:GE-C321-zurueckkehren}



\begin{quote}
Before the new one: say the German for to found, then the German for to come into being.
\end{quote}

\subsection*{You'll want to know: zurückkehren}
\textbf{zurückkehren} — \textquotedblleft{}to return, to come back\textquotedblright{}.

\begin{quote}
Nach zwanzig Jahren kehrte er in sein Dorf zurück.
\end{quote}

\begin{quote}
\emph{After twenty years he returned to his village.}
\end{quote}

Separable, with \emph{sein}: \emph{Er ist zurückgekehrt}. More formal than \emph{zurückkommen}.

\begin{cousinweb}
\emph{zurückkehren} joins \emph{zurück} \textquotedblleft{}back\textquotedblright{} and \emph{kehren} \textquotedblleft{}to turn\textquotedblright{}, from Old High German \emph{kēran}, Proto-Germanic \emph{karzijaną} (EtymDB): to turn back.
\end{cousinweb}

\begin{grammarlens}[title={What you've built}]
One more word for telling, explaining and describing at B1.
\end{grammarlens}

\subsection*{Your turn}
\begin{itemize}
\raggedright
  \item \emph{Say it:} \emph{zurückkehren}
  \item \emph{Say it:} \emph{zurückkehren}, once more
  \item \emph{Say it:} a sentence of your own with \emph{zurückkehren}
  \item \emph{Recall:} say the German for to found, then the German for to come into being, then say \emph{zurückkehren} again
\end{itemize}

\begin{tcolorbox}[breakable,colback=teal!4,colframe=teal!35!black,title={Before you move on}]
What does zurückkehren mean? (\textquotedblleft{}To return, to come back\textquotedblright{}.) Say it once more.
\end{tcolorbox}
